On the textbook maximisation-bias MDP our reimplementation reproduces the known behaviour: Q-learning overestimates the start-state value and takes the trap action far more often than Double Q-learning, and the effect grows with the number of actions at the trap state. On random 20-state MDPs the same comparison does not carry over at our budget. All four estimators underestimate the start-state value at the end of training, Q-learning the least, and Double Q-learning takes more suboptimal actions than Q-learning. Ablations are consistent with this being a finite-time effect of slower per-table learning and of greedy, adaptively collected data: with a warm start Double takes fewer suboptimal actions, with uniformly random behaviour the textbook sign pattern (Q-learning positive, Double negative) appears, and updating both tables on the same sample removes the benefit of Double entirely. Maxmin Q-learning with two estimators controls the bias at the price of action quality in the larger-action-space trap. Our evidence is limited to small tabular MDPs, one noise model and untuned shared hyper-parameters, so the practical message is cautious: judge a bias correction by the policy it yields and by the training regime, not by the sign of the bias alone.
